% =====================================================================
%  APÊNDICES -- Manual, seção 2.3.3
%  APÊNDICE = material ELABORADO PELO AUTOR (questionários, personas,
%  prompts, código-fonte). Pelo Manual, no sumário aparece só "APÊNDICES".
% =====================================================================
\begin{apendicesenv}
\partapendices
\addtocontents{toc}{\protect\setcounter{tocdepth}{-1}}

% ---------------------------------------------------------------------
\chapter{Questionário de avaliação das ideias}
\label{ap:questionario}

Este apêndice reproduz o questionário aplicado aos usuários reais, também
usado na geração das respostas sintéticas. O formulário foi feito no
Google Forms, com tempo estimado de resposta de 6 a 8 minutos, e abria com
o texto de apresentação e consentimento abaixo.

\begin{citacao}
Este questionário faz parte de um Trabalho de Conclusão de Curso em
Engenharia da Computação no Centro Universitário Ibmec. O objetivo é
entender o que motiva as pessoas a continuarem treinando lutas e artes
marciais, o que leva ao abandono da prática, e avaliar a receptividade a
um aplicativo de acompanhamento e engajamento voltado a academias.

As respostas são anônimas e serão utilizadas exclusivamente para fins
acadêmicos, de forma agregada, sem identificação individual. O contato
(quando fornecido) é guardado separadamente das respostas e usado só para
convite a testes. O fornecimento de contato ao final é opcional e serve
apenas para convite a testes futuros do aplicativo. Você pode desistir a
qualquer momento simplesmente fechando a página. O tratamento dos dados
segue a Lei Geral de Proteção de Dados (Lei 13.709/2018).

Em caso de dúvidas, entre em contato pelo e-mail
andreccoelho2004@gmail.com.
\end{citacao}

Depois das perguntas de perfil, o respondente era encaminhado a uma das
três seções conforme a resposta à quarta pergunta, e todos terminavam na
seção de encerramento. As alternativas das perguntas fechadas foram
reconstruídas a partir das respostas registradas no formulário e das
listas do \textit{script} de geração das respostas sintéticas, que
validava cada resposta contra as alternativas do formulário; por isso,
todas as alternativas listadas existem no formulário, mas alguma
alternativa nunca escolhida pode não aparecer.

{\small\SingleSpacing
\setlength{\columnsep}{0.8cm}
\begin{multicols}{2}
% Gerado a partir do cabecalho das planilhas de respostas, com as alternativas observadas.
\section*{Termo de consentimento e perfil (todos os respondentes)}
\begin{enumerate}[leftmargin=*, itemsep=1pt, parsep=0pt]
  \item Li as informações acima e concordo em participar da pesquisa.\\
        {\footnotesize\textit{Escolha única:} Sim, concordo em participar.}
  \item Faixa etária\\
        {\footnotesize\textit{Escolha única:} Menos que 18 anos; 18 a 24 anos; 25 a 34 anos; 35 a 44 anos; 45 a 54 anos; 55 anos ou mais.}
  \item Gênero\\
        {\footnotesize\textit{Escolha única:} Feminino; Masculino; Prefiro não dizer.}
  \item Qual das opções descreve melhor a sua relação com lutas ou artes marciais hoje?\\
        {\footnotesize\textit{Escolha única:} Treino atualmente; Já treinei, mas parei; Sou professor ou instrutor.}
\end{enumerate}

\section*{Seção do professor}
\begin{enumerate}[leftmargin=*, itemsep=1pt, parsep=0pt]
  \item Há quanto tempo você dá aulas de luta ou arte marcial?\\
        {\footnotesize\textit{Escolha única:} De 1 a 3 anos; De 3 a 5 anos; De 5 a 10 anos; Mais de 10 anos.}
  \item Quais modalidades você ensina?\\
        {\footnotesize\textit{Múltipla escolha:} Boxe; Muay Thai; Jiu-Jitsu; Karatê; Taekwondo; Judô; MMA.}
  \item Aproximadamente quantos alunos ativos você tem hoje? (na academia e particulares)\\
        {\footnotesize\textit{Escolha única:} Até 20; De 21 a 50; De 51 a 100; De 101 a 200; Mais de 200.}
  \item De cada 10 alunos que se matriculam, quantos ainda estão treinando depois de 6 meses?\\
        {\footnotesize\textit{Resposta numérica de 0 a 10.}}
  \item Em que momento a desistência costuma ser mais frequente?\\
        {\footnotesize\textit{Escolha única:} No primeiro mês; Entre 1 e 3 meses; Entre 3 e 6 meses; Não há um padrão claro.}
  \item Na sua percepção, quais são os principais motivos que levam os alunos a abandonarem o treino? (Selecione até 3 mais relevantes)\\
        {\footnotesize\textit{Múltipla escolha:} Falta de tempo ou conflito com a rotina; Questões financeiras; Não perceber evolução técnica; Ausência de metas claras; Lesão ou problema de saúde; Vergonha, insegurança ou intimidação; Falta de vínculo com a turma; Horários das aulas; Cansaço e desgaste da rotina de treino.}
  \item Como você acompanha a frequência dos seus alunos atualmente?\\
        {\footnotesize\textit{Escolha única:} Aplicativo de mensagens (WhatsApp); Caderno ou lista em papel; Sistema de gestão da academia; Planilha (Excel, Google Sheets); Não acompanho de forma sistemática.}
  \item Como funciona a graduação ou troca de faixa na sua academia, e como você controla quem está apto a graduar?\\
        {\footnotesize\textit{Resposta aberta.}}
  \item Quando um aluno some por algumas semanas, você costuma entrar em contato?\\
        {\footnotesize\textit{Escolha única:} Sim, sempre; Às vezes, quando percebo; Raramente.}
  \item Qual é hoje a sua maior dificuldade na gestão da turma?\\
        {\footnotesize\textit{Resposta aberta.}}
  \item Avalie de 1 (pouco útil) a 5 (muito útil) cada recurso:
  \begin{enumerate}[label=\alph*)]
    \item Pontos de experiência (XP) atribuídos automaticamente a cada treino
    \item Registro rápido de presença (check-in) a cada treino
    \item Contador de dias treinados e frequência ao longo das semanas
    \item Missões e desafios semanais criados por você para a turma
    \item Ranking de frequência entre os alunos da turma
    \item Medalhas e conquistas por marcos atingidos (100 treinos, 6 meses seguidos)
    \item Barra de progresso rumo à próxima graduação, visível para o aluno
    \item Alerta automático quando um aluno fica mais de 15 dias sem aparecer
    \item Relatório de frequência e evolução da turma
  \end{enumerate}
  \item Que riscos ou problemas você vê no uso de ranking e competição entre alunos dentro da academia?\\
        {\footnotesize\textit{Resposta aberta.}}
  \item Quanto tempo por semana você estaria disposto a dedicar para alimentar um aplicativo desse tipo?\\
        {\footnotesize\textit{Escolha única:} Menos de 5 minutos; De 5 a 15 minutos; De 15 a 30 minutos; Mais de 30 minutos.}
  \item Você utiliza algum sistema ou aplicativo de gestão hoje? Se sim, qual, e o que falta nele?\\
        {\footnotesize\textit{Resposta aberta.}}
  \item Se você pudesse pedir uma única funcionalidade para esse aplicativo, qual seria?\\
        {\footnotesize\textit{Resposta aberta.}}
  \item O que faria você desistir de usar o aplicativo depois de algumas semanas?\\
        {\footnotesize\textit{Resposta aberta.}}
  \item Existe alguma prática da sua academia que hoje funciona bem para segurar aluno e que poderia virar funcionalidade no aplicativo?\\
        {\footnotesize\textit{Resposta aberta.}}
\end{enumerate}

\section*{Seção do praticante}
\begin{enumerate}[leftmargin=*, itemsep=1pt, parsep=0pt]
  \item Qual modalidade você pratica atualmente?\\
        {\footnotesize\textit{Múltipla escolha:} Boxe; Muay Thai; Jiu-Jitsu; Karatê; Taekwondo; Judô; MMA.}
  \item Há quanto tempo você treina?\\
        {\footnotesize\textit{Escolha única:} Menos de 3 meses; De 3 a 6 meses; De 6 meses a 1 ano; De 1 a 3 anos; Mais de 3 anos.}
  \item Quantas vezes por semana você treina, em média?\\
        {\footnotesize\textit{Escolha única:} 1 vez; 2 vezes; 3 vezes; 4 vezes; 5 vezes ou mais.}
  \item O que faz você continuar treinando? (Selecione até 3)\\
        {\footnotesize\textit{Múltipla escolha:} Saúde física e condicionamento; Saúde mental e alívio do estresse; Amizades e convívio com a turma; Defesa pessoal; Disciplina e rotina; Competição; Evolução técnica e aprendizado; Estética e composição corporal; Relação com o professor.}
  \item De 1 a 5, o quanto você consegue perceber a sua própria evolução técnica?\\
        {\footnotesize\textit{Escala de 1 a 5.}}
  \item Você sabe exatamente o que precisa fazer para chegar à próxima graduação ou faixa?\\
        {\footnotesize\textit{Escolha única:} Sim, sei com clareza; Não faço ideia; Tenho uma ideia aproximada; Minha modalidade ou academia não trabalha com graduação.}
  \item Você já pensou em parar de treinar em algum momento? O que quase o/a fez você parar?\\
        {\footnotesize\textit{Resposta aberta.}}
  \item Você já parou de treinar e voltou depois?\\
        {\footnotesize\textit{Escolha única:} Sim; Não.}
  \item Você acompanha a sua frequência de treinos de alguma forma?\\
        {\footnotesize\textit{Escolha única:} Só mentalmente, sem registro; Não acompanho; Sim, em anotações próprias ou planilha; Sim, em aplicativo.}
  \item Avalie de 1 (pouco útil) a 5 (muito útil) cada recurso:
  \begin{enumerate}[label=\alph*)]
    \item Pontos de experiência (XP) ganhos a cada treino
    \item Registro rápido de presença (check-in) a cada treino
    \item Contador de dias treinados e frequência ao longo das semanas
    \item Sequência de dias ou semanas treinadas sem quebrar (ofensiva)
    \item Missões e desafios semanais propostos pelo professor
    \item Ranking de frequência entre os alunos da turma
    \item Medalhas por marcos atingidos (100 treinos, 1 ano de prática)
    \item Barra de progresso mostrando o quanto falta para a próxima graduação
    \item Histórico completo de presença e evolução pessoal
    \item Mural ou feed da turma, com conquistas dos colegas
    \item Lembretes automáticos nos dias em que você costuma treinar
  \end{enumerate}
  \item De 1 a 5, o quanto você usaria um aplicativo assim se os seus dados fossem visíveis apenas para você e para o seu professor?\\
        {\footnotesize\textit{Escala de 1 a 5.}}
  \item Você já usou algum aplicativo com pontos, metas ou sequências que ajudou a manter um hábito? Qual, e o que funcionou nele?\\
        {\footnotesize\textit{Resposta aberta.}}
  \item Se esse aplicativo existisse, que funcionalidade você gostaria que ele tivesse?\\
        {\footnotesize\textit{Resposta aberta.}}
  \item O que faria você desinstalar o aplicativo?\\
        {\footnotesize\textit{Resposta aberta.}}
  \item Que informação sobre o seu treino você gostaria de ver registrada e que hoje se perde?\\
        {\footnotesize\textit{Resposta aberta.}}
\end{enumerate}

\section*{Seção do ex-praticante}
\begin{enumerate}[leftmargin=*, itemsep=1pt, parsep=0pt]
  \item Qual modalidade você praticava?\\
        {\footnotesize\textit{Escolha única:} Boxe; Muay Thai; Jiu-Jitsu; Karatê; Taekwondo; Judô; MMA.}
  \item Por quanto tempo você treinou antes de parar?\\
        {\footnotesize\textit{Escolha única:} Menos de 1 mês; De 1 a 3 meses; De 3 a 6 meses; De 6 meses a 1 ano; De 1 a 3 anos; Mais de 3 anos.}
  \item Há quanto tempo você parou?\\
        {\footnotesize\textit{Escolha única:} Menos de 6 meses; De 6 meses a 1 ano; De 1 a 3 anos; Mais de 3 anos.}
  \item Quantas vezes por semana você treinava, em média?\\
        {\footnotesize\textit{Escolha única:} 1 vez; 2 vezes; 3 vezes; 4 vezes; 5 vezes ou mais.}
  \item Quais fatores contribuíram para você parar de treinar? (Selecione até 3)\\
        {\footnotesize\textit{Múltipla escolha:} Falta de tempo ou conflito com a rotina; Questões financeiras; Não perceber evolução técnica; Ausência de metas claras; Lesão ou problema de saúde; Vergonha, insegurança ou intimidação; Falta de vínculo com a turma; Pouca proximidade com o professor; Horários das aulas; Mudança de cidade ou bairro; Cansaço e desgaste da rotina de treino; Perda de interesse pela modalidade.}
  \item Entre os fatores que você marcou, qual foi o principal?\\
        {\footnotesize\textit{Escolha única:} Falta de tempo ou conflito com a rotina; Questões financeiras; Não perceber evolução técnica; Ausência de metas claras; Lesão ou problema de saúde; Vergonha, insegurança ou intimidação; Falta de vínculo com a turma; Pouca proximidade com o professor; Horários das aulas; Mudança de cidade ou bairro; Cansaço e desgaste da rotina de treino; Perda de interesse pela modalidade.}
  \item Conte com suas palavras como foi o processo até você parar de treinar\\
        {\footnotesize\textit{Resposta aberta.}}
  \item De 1 a 5, o quanto você sentia que estava evoluindo tecnicamente pouco antes de parar?\\
        {\footnotesize\textit{Escala de 1 a 5.}}
  \item Você sabia o que precisava fazer para chegar à próxima graduação?\\
        {\footnotesize\textit{Escolha única:} Sim, sabia com clareza; Tinha uma ideia aproximada; Não fazia ideia.}
  \item Depois que você parou de aparecer, alguém da academia tentou te chamar de volta?\\
        {\footnotesize\textit{Escolha única:} Sim, o professor ou instrutor me chamou de volta; Ninguém entrou em contato; Sim, a administração ou secretaria entrou em contato; Alguém da academia falou comigo, mas apenas socialmente, sem me chamar para voltar.}
  \item Se alguém entrou em contato, isso fez você considerar voltar?\\
        {\footnotesize\textit{Escolha única:} Sim, considerei, mas não voltei; Não fez diferença; Não se aplica, ninguém me chamou de volta; Sim, cheguei a voltar.}
  \item Você tem vontade de voltar a treinar?\\
        {\footnotesize\textit{Escolha única:} Sim, pretendo voltar em breve; Talvez, um dia; Não pretendo voltar.}
  \item O que faria você voltar a treinar hoje?\\
        {\footnotesize\textit{Resposta aberta.}}
  \item Avalie de 1 (pouco útil) a 5 (muito útil) cada recurso:
  \begin{enumerate}[label=\alph*)]
    \item Pontos de experiência (XP) ganhos a cada treino
    \item Registro rápido de presença (check-in) a cada treino
    \item Contador de dias treinados e frequência ao longo das semanas
    \item Sequência de dias ou semanas treinadas sem quebrar (ofensiva)
    \item Missões e desafios semanais propostos pelo professor
    \item Ranking de frequência entre os alunos da turma
    \item Medalhas por marcos atingidos
    \item Barra de progresso mostrando o quanto falta para a próxima graduação
    \item Histórico completo de presença e evolução pessoal
    \item Mural ou feed da turma, com conquistas dos colegas
    \item Mensagem automática do professor após alguns dias de ausência
    \item Lembretes automáticos nos dias de treino
  \end{enumerate}
  \item Sendo sincero, você acredita que um aplicativo assim teria mudado a sua decisão de parar?\\
        {\footnotesize\textit{Escolha única:} Não, o motivo pelo qual parei não tem relação com isso; Talvez, teria ajudado, mas não sei se seria suficiente; Sim, provavelmente eu teria continuado.}
  \item Algum desses recursos teria te incomodado ou te afastado?\\
        {\footnotesize\textit{Escolha única:} Nenhum me incomodaria; Lembretes e notificações; Mural com conquistas dos colegas; Ranking visível para outros alunos; Mensagem do professor cobrando presença.}
  \item Que funcionalidade um aplicativo da academia precisaria ter para fazer você voltar e permanecer?\\
        {\footnotesize\textit{Resposta aberta.}}
  \item O que faria você desinstalar um aplicativo desse tipo?\\
        {\footnotesize\textit{Resposta aberta.}}
\end{enumerate}

\section*{Encerramento (todos os respondentes)}
\begin{enumerate}[leftmargin=*, itemsep=1pt, parsep=0pt]
  \item Deseja acrescentar algum comentário sobre o tema?\\
        {\footnotesize\textit{Resposta aberta.}}
  \item Você aceitaria participar de um teste do aplicativo ou de uma conversa rápida de acompanhamento?\\
        {\footnotesize\textit{Escolha única:} Sim; Não.}
  \item Deixe um e-mail ou telefone para contato (opcional).\\
        {\footnotesize\textit{Resposta aberta.}}
\end{enumerate}


\end{multicols}
}

% ---------------------------------------------------------------------
\chapter{Artefatos técnicos no repositório}
\label{ap:artefatos}

O código-fonte e a documentação técnica do RoundUp estão no repositório
público do projeto, na versão entregue com este relatório, marcada com a
\textit{tag} \texttt{relatorio-v2}:
\url{https://github.com/andreccoelho/roundup_app/tree/relatorio-v2}.
O \autoref{qua:artefatos} indica onde está cada artefato; os caminhos são
relativos à raiz do repositório.

\begin{quadro}[!htb]
  \caption{Artefatos técnicos do projeto no repositório}
  \label{qua:artefatos}
  \centering
  \footnotesize
  \begin{tabular}{|p{4.6cm}|p{5.2cm}|p{3.4cm}|}
    \hline
    \textbf{Artefato} & \textbf{Caminho} & \textbf{Relação com o texto} \\ \hline
    Lista oficial de requisitos (RF, RN e RNF) & \path{docs/requisitos/requisitos-v1-2026-09.csv} & \autoref{cap:requisitos} \\ \hline
    Especificação dos 30 casos de uso & \path{docs/requisitos/especificacao-casos-uso.md} & \autoref{sec:modelagem-funcional} \\ \hline
    Fontes dos diagramas de casos de uso (PlantUML) & \path{docs/relatorio-v2/figuras/uml-fonte/} & \autoref{cap:modelagem} \\ \hline
    Registro das decisões técnicas & \path{docs/decisoes-tecnicas.md} & \autoref{cap:implementacao} \\ \hline
    Código do aplicativo (React Native e Expo) & \path{src/} & \autoref{cap:implementacao} \\ \hline
    Funções de servidor (Cloud Functions) & \path{functions/src/} & \autoref{cap:implementacao} \\ \hline
    Regras de segurança do Firestore e do Storage & \path{firestore.rules}, \path{storage.rules} & \autoref{cap:modelagem} e \autoref{cap:implementacao} \\ \hline
        Testes automatizados (regras, funções e utilitários) & \path{tests/}, \path{src/utils/*.test.ts} & \autoref{cap:testes} \\ \hline
    Roteiro de validação das notificações em aparelho & \path{docs/ROTEIRO_TESTE_NOTIFICACOES.md} & \autoref{cap:implementacao} e \autoref{cap:testes} \\ \hline
    \textit{Scripts} de migração de dados (medalhas, faixas, público das missões e professor atribuído) & \path{scripts/migrar-*.cjs} & \autoref{cap:implementacao} \\ \hline
    Fluxos de publicação e preparação de testes (GitHub Actions) & \path{.github/workflows/} & \autoref{cap:implementacao} \\ \hline
    \textit{Script} das contas de demonstração (academia, professor e aluno) & \path{scripts/seed-demo.cjs} & -- \\ \hline
    Dados anonimizados e \textit{scripts} de análise do questionário & \path{docs/relatorio-v2/dados/}, \path{docs/relatorio-v2/scripts/} & \autoref{cap:testes} \\ \hline
    Fichas das personas sintéticas & \path{docs/pesquisa/personas-sinteticas.md} & \autoref{sec:val-personas} \\ \hline
    Geração e envio das respostas sintéticas & \path{docs/pesquisa/respostas-sinteticas/} & \autoref{sec:val-personas} \\ \hline
  \end{tabular}
  \fonte{Elaborado pelo autor (2026).}
\end{quadro}

\end{apendicesenv}
